\subsection{The augmentation changes the conditional risk target}
\label{sec:noise-augmentation}
The training code perturbs observed positions by $\xi_k$ and uses the adjusted
next-position target $x_{t+1}+\xi_t$. Its acceleration target is therefore
\begin{align}
Y_t
&=(x_{t+1}+\xi_t)-2(x_t+\xi_t) \nonumber\\
&\quad +(x_{t-1}+\xi_{t-1}) \nonumber\\
&=A_t-\eta_t,\qquad \eta_t=\xi_t-\xi_{t-1},
\label{eq:noise-target-identity}
\end{align}
where $A_t=x_{t+1}-2x_t+x_{t-1}$ is clean acceleration in stored-frame
displacement units. For a predictor evaluated on noisy history $\widetilde H$,
write $R=A_t-\mu(\widetilde H)$. Then
\begin{align}
\mathbb E[\|Y_t-\mu\|^2\mid\widetilde H]
={}&\mathbb E[\|R\|^2\mid\widetilde H] \nonumber\\
&+\mathbb E[\|\eta_t\|^2\mid\widetilde H]
-2\mathbb E[R^\top\eta_t\mid\widetilde H].
\label{eq:noise-risk-decomposition}
\end{align}
Applying the saved affine normalization gives the same identity with scaled
residual and noise terms. Independence of injected noise and clean trajectories
before conditioning does not justify discarding the cross term: the observed
history and predictor depend on the noise. Calibration for this augmented
regression task consequently does not by itself establish clean-evaluation
residual calibration.
This does not invalidate adding the marginal noise variance to the training
normalizer: if $A_t$ and $\eta_t$ are independent before conditioning, their
unconditional variances add. Marginal normalization and conditional residual
calibration are different requirements.

For a scalar illustration in unnormalized units, let clean acceleration be zero
and $V\sim\mathcal N(0,\tau^2)$. The five observed velocities of a six-frame
history are $U_k=V+\sum_{j=1}^kz_j$, with independent
$z_j\sim\mathcal N(0,\delta)$, $\delta=s^2/5$, and target
$Y=-\sum_{j=1}^5z_j=V-U_5$. Because $U_k-U_{k-1}=z_k$ for $k\ge2$, only
$U_1=V+z_1$ informs the posterior of $V$. For $\tau,s>0$, Gaussian conditioning
therefore gives
\begin{align}
\mu(U)&=aU_1-U_5,& q(U)&=h,\nonumber\\
a&=\frac{\tau^2}{\tau^2+\delta},&
h&=\frac{\tau^2\delta}{\tau^2+\delta}.
\label{eq:noise-toy-posterior}
\end{align}
These are the exact augmented-target conditional mean and variance. At a
noiseless history $U_k=v$, the continuous Gaussian Bayes mean is $-(1-a)v$,
although clean acceleration is zero. Its clean-population MSE is
\begin{equation}
\mathbb E_V[(0-\mu(V,\ldots,V))^2]
=(1-a)^2\tau^2=(1-a)h,
\label{eq:noise-toy-clean-risk}
\end{equation}
while the head remains $h$. This aggregate mismatch need not hold in the same
direction for every velocity. Exact clean histories form a measure-zero
diagonal under the noisy training law; the statement uses the continuous Bayes
version, not arbitrary measurable minimizers. The example establishes a
target-and-input distinction, not the magnitude or direction of the effect in
our particle experiments. It does not imply that augmentation harms simulation.
